\subsection{Adaptive Virtual Receiver Generation}
\label{subsec:adaptive-virtual-rx-generation}

The coherence cache described in Section~\ref{sec:sionna-cache-mechanism} is reactive: a new ray tracing call is issued only after a cache entry becomes invalid.
Adaptive virtual receiver generation is a complementary proactive mechanism that reduces the number of such cache misses.
During each cache refresh, a small number of temporary receiver positions are placed ahead of each UE along its trajectory.
These positions are not new network nodes; they are spatial samples for the same UE at predicted future locations.
Including them in the same Sionna RT call returns both present and near-future channel samples at no additional ray tracing overhead.

Let the UE position and speed be
\begin{equation}
    \mathbf{x}_r(t)=[x_r(t),y_r(t),z_r(t)]^{\mathsf{T}},
    \qquad v_r .
\end{equation}
Prediction of virtual receiver positions is not performed for a stationary receiver (i.e., $v_r \leq 10^{-6}$~m/s).
For a moving UE, the distance to the nearest transmitter is
\begin{equation}
    d_r(t)=\min_{b\in\mathcal{B}}
    \left\|\mathbf{x}_r(t)-\mathbf{x}_b\right\|_2 .
\end{equation}
The valid displacement $\Delta_r$ is computed using the coherence formula
from Section~\ref{subsec:cache-validity}, based on $d_r(t)$ and a default
$\Delta_{\min}=0.5$ m.
The current sample is predicted to expire after
\begin{equation}
    T_{\mathrm{expire},r}=\frac{\Delta_r}{v_r}.
\end{equation}
This is the time for the UE to travel the full valid displacement at its current speed.
Virtual receivers are generated only if $T_{\mathrm{expire},r}<T_{\min}$, with $T_{\min}=1$ s.

The generation of virtual receivers is closely tied to the underlying mobility model of the UE.
For UEs with predictable trajectories, such as those using a Waypoint mobility model, the direction of motion is stable and virtual receivers are reliably placed along the projected path.
In contrast, for unpredictable movement patterns, such as a random walk, predicting future positions is unreliable.

To differentiate these scenarios, the framework assesses trajectory stability by checking that the dot product of the current and previous unit direction vectors satisfies
\begin{equation}
    \mathbf{u}_r(t-\Delta t)^{\mathsf{T}}\mathbf{u}_r(t)\geq \rho_{\min},
\end{equation}
where $\rho_{\min}=0.7$, and the prediction direction is defined by the normalized horizontal displacement,
\begin{equation}
    \mathbf{u}_r(t)=
    \frac{[\Delta x_r(t),\Delta y_r(t),0]^{\mathsf{T}}}
    {\sqrt{\Delta x_r(t)^2+\Delta y_r(t)^2}} .
\end{equation}
If this stability condition is not met, no virtual receivers are generated.
This prevents wasting computational resources on ray tracing for locations the UE is unlikely to visit.

For stable trajectories, the inter-sample spacing $s_r$ and prediction horizon $D_r$ are
\begin{equation}
    s_r=\beta\Delta_r, \qquad D_r=v_rT_h,
\end{equation}
with $\beta=1.1$ and $T_h=3$ s.
The number of virtual receivers is
\begin{equation}
    K_r=\min\left(K_{\max},\left\lceil\frac{D_r}{s_r}\right\rceil\right),
\end{equation}
where $K_{\max}=3$.
Their positions are calculated as
\begin{equation}
    \mathbf{x}^{\mathrm{virt}}_{r,k}(t)
    =
    \mathbf{x}_r(t)+ks_r\mathbf{u}_r(t),
    \qquad k=1,\ldots,K_r .
\end{equation}
Thus, only positions along the UE trajectory are sampled, and the receiver height remains unchanged.